\section{襄国、成都与建康之间：三三〇至三四九年的多重重组}
\label{sec:eastern-jin-post-zhao-chenghan}

三二九年前赵灭亡后，石勒在次年于襄国称帝，后赵成为北方最强的政权之一。称帝与改元可靠，官制仿效到什么程度、羯胡与汉人官僚如何共处、地方命令是否真正执行到各城镇，不能由《晋书》〈石勒载记〉的帝国叙事直接推出。后赵能调动河北的资源，也意味着黄河流域的城郭、军队与家庭必须承担新的征发与秩序。

\subsection{继承危机没有只发生在建康}

三三三年石勒死，太子石弘继位，石虎掌握军政；次年石虎废杀石弘，自立为天王。遗诏、石弘的自主空间和石虎取得支持的过程都有争议，废立与权力转移却无可怀疑。后赵的高压扩张暂时得到更集中的指挥，也把继承危机压入宗室、军队和宫廷之间。三四九年石虎死后，诸子与养孙围绕邺城、军队和皇位相争，说明先前的集中并没有建立能够自我约束的继承程序。

西南的成汉走过相似而不相同的道路。李雄三三四年去世，李班继位后旋即被李期夺权；三三八年李寿再废李期，改国号为汉。成都宫廷中谁支持哪一位宗室、各次废立的罪状是否可靠，不能全凭后出的载记裁断；频繁改号与政变确实削弱了政权的正当性和应变能力。成汉并非因“偏远”而天然脆弱，它需要在巴蜀的官僚、道路和军队之间维持协调，而继承冲突正在消耗这种协调。

\subsection{辽西、丸都与区域的道路}

三三五年慕容皝在辽西称燕王，次年击败段辽并整合部众；三三八年石虎进攻前燕未能摧毁其核心地区。战场细节、部众数量和所谓降附者的迁徙路线有争议，但辽西通道、棘城周边地形和可动员的人口，使前燕得以在后赵的压力下继续发展。三四二年慕容皝进攻高句丽都城丸都，远征和城邑受创可由《晋书》、后出的《三国史记》及《资治通鉴》交叉辨认；俘获人数、王室损失与双方叙事的差异不应被抹平。

这些北方和东北的战事并不是建康能够直接指挥的边缘故事。它们不断改变流民、使节、军镇和地方集团可以选择的道路，也决定东晋北伐面对的是哪一种政权组合。东晋的北方目标因而不能只写成恢复洛阳的愿望，还要看见每一条从江淮通往关中、河北和辽东的路径，都由不同的城郭和人际网络控制。

\subsection{墓志与巴蜀的再接合}

三四一年王兴之夫妇墓志下葬，南京象山的发掘为建康周边侨姓官员家庭、婚姻和官职提供了带纪年的物证。墓志可以校正姓名、日期和某些仕宦关系，却是家族的公开自我呈现；追赠官、世系和道德语言不能替代更广泛的社会统计。它的价值正在于提示我们，东晋的政治并不只存在于帝纪中的王、桓等人，也存在于迁徙家庭如何安置、婚配和书写身份的地方网络。

三四六年桓温自荆州西征成汉，次年攻入成都，李势降，东晋恢复益州控制。出师、主要路线与成都降服可靠，兵数、补给、战后安置及蜀地整合的速度均有争议。东晋取得上游的税源和战略纵深，也显著抬高了桓温的个人声望；巴蜀重新进入建康的任命网络，却不等于当地社会立即按江南中枢的意图运转。成汉覆灭与后赵继承危机相互并行，使三四九年以后的北方又一次进入更剧烈的重组。

\subsection{材料的边界}

《晋书》〈石勒载记〉〈石虎载记〉、慕容与李氏载记，《华阳国志》以及《资治通鉴》共同保存这一时期的政治主线；十六国史多已散佚，后代辑佚文字能补国号和地名，却不能被当作完整原史。王兴之夫妇墓志和发掘报告提供独立的地域与纪年锚点，但不能代表建康的全部人口。由此可见的，是襄国、成都、棘城和建康之间并无一条单向的强弱线，而是多种继承、迁徙和军政网络同时改变的历史空间。


% 年序补线：本节将既有账本中尚未初次落入正文的条目接入相应历史场景。
\subsection{北方重组：邺城、关中与河西的去向}
把\eventanchor{JH-0341-WANG-XINGZHI-EPITAPH}{王兴之夫妇墓志下葬，为建康周边侨姓官员家庭、婚姻与官职提供有纪年实物材料}放回咸康七年（341年），可见东晋江南正围绕王兴之夫妇墓志下葬，为建康周边侨姓官员家庭、婚姻与官职提供有纪年实物材料作出处理。 可见的选择边界来自南渡士族在建康周边定居并通过官职、婚姻和墓葬维系社会身份，并留下出土墓志为仅依正史重建东晋精英网络提供独立的日期和地域锚点。 这段叙述只越不过《晋书·成帝纪》；〈王兴之夫妇墓志〉；南京象山王兴之夫妇墓发掘报告所能承载的范围，墓志纪年与墓主姓名可靠；官衔追赠、家世叙述和墓志能否代表更广泛社会有争议。

把\eventanchor{JH-0342-MURONG-HUANG-GOGURYEO}{慕容皝进攻高句丽都城丸都，掳掠人口并改变辽东力量关系}放回咸康八年（342年），可见前燕与高句丽正围绕慕容皝进攻高句丽都城丸都，掳掠人口并改变辽东力量关系作出处理。 这里的资源与选择关系可从前燕为控制辽东通道并消除后方威胁，选择对高句丽核心地区发动远征看出，结果使高句丽短期受挫而前燕获得东北声望，双方冲突成为区域政治的持续因素。 以《晋书·慕容皝载记》；《三国史记·高句丽本纪》；《资治通鉴·晋纪》交叉检验，能够确认的止于此；远征与丸都受创可靠；路线、俘获人数、王室损失和双方叙事差异有争议。

永和二年（346年）的东晋与成汉并非抽象背景；\eventanchor{JH-0346-HUAN-WEN-INVADES-CHENGHAN}{桓温受命自荆州西征成汉，沿长江上游进入巴蜀}对应的是桓温受命自荆州西征成汉，沿长江上游进入巴蜀。 可见的选择边界来自成汉内乱削弱成都防御，东晋希望以西征提升政权威望并控制上游，并留下东晋军队突破峡江防线，巴蜀政权面临终局性军事压力。 就材料边界说，《晋书·桓温传》；《晋书·穆帝纪》；《华阳国志·李特雄期寿势志》；《资治通鉴·晋纪》支持这一轮廓；出师、主要路线和桓温统帅地位可靠；兵数、补给和朝廷授命范围有争议。

永和三年（347年），东晋与成汉的\eventanchor{JH-0347-CHENGHAN-FALL}{桓温攻入成都，李势投降，成汉灭亡，东晋恢复益州控制}使桓温攻入成都，李势投降，成汉灭亡，东晋恢复益州控制成为可追索的行动。 它把成汉军政因继承冲突而弱，桓温军利用长江水陆通道直逼成都带入新的局面，令东晋重获巴蜀税源与上游战略纵深，桓温的个人威望显著上升成为后来人的问题。 以《晋书·穆帝纪》；《晋书·桓温传》；《晋书·李势载记》；《华阳国志·李特雄期寿势志》交叉检验，能够确认的止于此；成都降服与成汉灭亡可靠；李势投降条件、战后安置和东晋对蜀地的实际整合有争议。

把\eventanchor{JH-0349-SHI-HU-DEATH}{石虎去世，诸子与养孙冉闵围绕都城邺城和军队展开继承争夺}放回永宁元年（349年），可见后赵正围绕石虎去世，诸子与养孙冉闵围绕都城邺城和军队展开继承争夺作出处理。 这一节点将石虎长期以清洗维系统治，未能形成稳定的继承与军权交接机制与后赵中央迅速分裂，华北各地军队和地方集团获得自立空间接合，却不预设两者之间只有一种因果。 《晋书·石虎载记》；《晋书·冉闵载记》；《资治通鉴·晋纪》使这项处置可被追查，不足以抹平石虎卒与继承危机可靠；遗诏、石世即位程序和冉闵早期角色有争议。

\subsection{前燕与东晋：北伐并未恢复北方秩序}
永宁二年（350年），前秦前身的\eventanchor{JH-0350-FU-HONG-CLAIMS-QIN}{氐族首领苻洪在关中东部聚众，自称三秦王，前秦政权的军事基础开始形成}写下氐族首领苻洪在关中东部聚众，自称三秦王，前秦政权的军事基础开始形成这一处具体变化。 行动者受后赵崩解使关中东部失去稳定统治，苻洪部可吸纳流民和旧军人限制，随后苻氏取得建国基础，关中随后成为前秦与其他政权竞争的核心。 所据《晋书·苻洪载记》；《晋书·苻健载记》；《资治通鉴·晋纪》留下可核的线索，同时提醒读者苻洪称号与聚众可靠；部众来源、与冉魏和东晋的关系及控制区域有争议。

在永宁二年（350年），\eventanchor{JH-0350-RAN-MIN-WEI}{冉闵废石祗后在邺城称帝，建立冉魏并在内战中针对羯胡及后赵旧部实施暴力清洗}把冉魏与后赵遗众置于冉闵废石祗后在邺城称帝，建立冉魏并在内战中针对羯胡及后赵旧部实施暴力清洗的历史现场。 它不是孤立转折：后赵继承战使冉闵掌握邺城军队，他以排斥羯胡和恢复魏号争取部分支持构成背景，暴力加剧人口逃散和政治碎片化，冉魏同时遭前燕、前秦等多方压力则把压力送往下一阶段。 材料来自《晋书·冉闵载记》；《资治通鉴·晋纪》；《十六国春秋》辑佚；因此称帝与后赵旧部受大规模杀害可靠；受害群体界定、死亡数字和各地暴力的因果关系有争议。

永宁三年（351年）的年序到\eventanchor{JH-0351-FU-JIAN-FOUNDS-FORMER-QIN}{苻健在长安称天王，建立前秦并开始重建关中官署}时，前秦须面对苻健在长安称天王，建立前秦并开始重建关中官署。 它不是孤立转折：苻洪死后苻健继承部众，后赵与冉魏无力稳定关中构成背景，前秦以长安为基地进入十六国竞争，并逐步吸纳关中人口与官僚则把压力送往下一阶段。 据《晋书·苻健载记》；《资治通鉴·晋纪》；《十六国春秋》辑佚，可确认的是基本顺序和所列行动；称天王、入长安和前秦建国可靠；官制规模、城内接收过程与地方归附程度有争议。

\eventanchor{JH-0352-RAN-MIN-DEFEATED}{冉闵在魏昌一带败于前燕，被俘后处死，冉魏政权迅速瓦解}把永和八年（352年）的冉魏与前燕与冉闵在魏昌一带败于前燕，被俘后处死，冉魏政权迅速瓦解这一材料所见的处置连在一起。 其代价落在冉魏多面受敌且粮饷不足，前燕利用骑兵和河北联盟发动进攻所限定的处境中，并以前燕进入中原政治竞争，后赵崩解后华北并未恢复稳定统治延续下去。 就材料边界说，《晋书·冉闵载记》；《晋书·慕容俊载记》；《资治通鉴·晋纪》支持这一轮廓；冉闵被俘处死和冉魏衰亡可靠；战场细节、俘后处置和民间支持范围有争议。

在永和九年三月三日（353年）这一个节点，东晋通过\eventanchor{JH-0353-LANTING-GATHERING}{王羲之与会稽士人在兰亭修禊并作序，成为东晋士人交游与书写文化的标志性事件}显出王羲之与会稽士人在兰亭修禊并作序，成为东晋士人交游与书写文化的标志性事件的压力。 前因和后效须分开读：前者是会稽士族在相对安定的江南组织礼俗与文会，书写成为维系名望网络的媒介，后者是兰亭文本在后世反复摹写和政治化，亦提示东晋文化史不能只由宫廷与战事叙述。 材料的可说范围由《晋书·王羲之传》；《世说新语》及刘孝标注；《兰亭序》传世摹本与刻石材料界定；修禊和序文形成大体可靠；与会者名单、作品原貌和传世摹本关系有争议。

在永和十年（354年），\eventanchor{JH-0354-HUAN-WEN-NORTHERN-EXPEDITION}{桓温自荆州北伐前秦，进至关中外围后因粮运困难撤退}把东晋与前秦置于桓温自荆州北伐前秦，进至关中外围后因粮运困难撤退的历史现场。 这一节点将收复蜀地后桓温声望上升，东晋希望借前秦未稳扩大北方影响与补给线限制显示东晋难以长期经营关中，桓温的北伐成为军阀政治资源接合，却不预设两者之间只有一种因果。 这段叙述只越不过《晋书·桓温传》；《晋书·穆帝纪》；《晋书·苻健载记》；《资治通鉴·晋纪》所能承载的范围，北伐、进军关中和撤军可靠；军队数量、长安周边战况与撤退责任有争议。

沿着永和十一年（355年），\eventanchor{JH-0355-LATER-ZHAO-ENDS}{石祗在襄国被杀，后赵最后的政治中心消失}让后赵残余与冉魏的石祗在襄国被杀，后赵最后的政治中心消失进入连续叙事。 它不是孤立转折：石虎死后的继承战与冉魏崛起不断削弱襄国，石祗无法重建有效联盟构成背景，后赵正式退出竞争，前燕、前秦和各地军镇争夺其遗留人口与城镇则把压力送往下一阶段。 材料的可说范围由《晋书·石虎载记下》；《晋书·冉闵载记》；《资治通鉴·晋纪》界定；石祗被杀与后赵终结大体可靠；行凶者动机、残余军队去向和终结时间界定有争议。

把\eventanchor{JH-0356-HUAN-WEN-RECOVERS-LUOYANG}{桓温北伐进入洛阳，修复晋帝陵并短暂恢复东晋在中原的象征性存在}放回永和十二年（356年），可见东晋与前燕正围绕桓温北伐进入洛阳，修复晋帝陵并短暂恢复东晋在中原的象征性存在作出处理。 这里的资源与选择关系可从后赵瓦解提供北伐窗口，桓温以恢复晋室旧都与陵寝凝聚政治声望看出，结果使东晋未能建立持久的河南统治，但洛阳成为南北争夺的正统象征。 证据不允许把场景填满：《晋书·桓温传》；《晋书·穆帝纪》；《资治通鉴·晋纪》仅能支持核心事项，入洛阳与修陵大体可靠；控制范围、停留时间和地方响应程度有争议。



\subsection{360—369年：材料所见的连续压力}
后赵瓦解后的城邑、河谷和旧军镇成为新政权争取的节点；东晋北伐与江南守备亦受同一套水陆条件限制。

把公元360年（年序桥接）放回连续叙事，\eventanchor{JH-DENS-360-A}{前燕、前秦与东晋围绕河北、关中和淮汉走廊轮番调兵，扩张依赖州郡供给与地方合作}保留前燕、前秦与东晋围绕河北、关中和淮汉走廊轮番调兵，扩张依赖州郡供给与地方合作，\eventanchor{JH-DENS-360-B}{东晋、前燕与前秦的朝廷和州郡围绕户籍、田租、军粮与转运的安排进入材料}让东晋、前燕与前秦的朝廷和州郡围绕户籍、田租、军粮与转运的安排进入材料继续约束东晋、前燕与前秦的行动。 它们不是由年份自动推出的因果，而是让前后事件之间的资源、道路、户籍或地方网络保持可见。 就可说范围而言，前一判断止于《晋书》哀海西简文诸帝纪；《晋书》载记；《资治通鉴·晋纪》，本纪与编年材料主要确证政权年序和精英行动；对基层执行、疆界和人口数量的沉默不得以叙事补足；后一判断止于《水经注》；墓志、城址与十六国史辑佚，志书、诏令与本纪可显示制度设计或战争需求，不能直接换算赋额、仓储或实际负担。

公元361年（年序桥接），\eventanchor{JH-DENS-361-A}{前燕、前秦与东晋围绕河北、关中和淮汉走廊轮番调兵，扩张依赖州郡供给与地方合作}和\eventanchor{JH-DENS-361-B}{东晋、前燕与前秦的流民、侨置户、军镇居民或地方家族的安置与编户问题构成社会背景}共同避免把东晋、前燕与前秦压缩为单一都城事件：前者关注前燕、前秦与东晋围绕河北、关中和淮汉走廊轮番调兵，扩张依赖州郡供给与地方合作，后者关注东晋、前燕与前秦的流民、侨置户、军镇居民或地方家族的安置与编户问题构成社会背景。 如此处理不是用空白填满史实，而是避免后来的转折抹去本年仍在运作的条件。 《晋书》哀海西简文诸帝纪；《晋书》载记；《资治通鉴·晋纪》使前一线索可追，本纪与编年材料主要确证政权年序和精英行动；对基层执行、疆界和人口数量的沉默不得以叙事补足；《水经注》；墓志、城址与十六国史辑佚限制后一线索的说法，墓志、家族文本和侨置名目各自有选择性，不能据此统计迁徙规模或概括整体社会态度。

在公元362年（年序桥接）的并行行动者之间，\eventanchor{JH-DENS-362-A}{前燕、前秦与东晋围绕河北、关中和淮汉走廊轮番调兵，扩张依赖州郡供给与地方合作}将前燕、前秦与东晋围绕河北、关中和淮汉走廊轮番调兵，扩张依赖州郡供给与地方合作落到时间位置，\eventanchor{JH-DENS-362-B}{东晋、前燕与前秦的寺院、僧侣、道士和施主网络在材料中连接都城、州郡与边地}从东晋、前燕与前秦的寺院、僧侣、道士和施主网络在材料中连接都城、州郡与边地补足资源与空间的视角。 这里保留的是可核的连续性，而非对基层反应、疆界或人数的无据补写。 《晋书》哀海西简文诸帝纪；《晋书》载记；《资治通鉴·晋纪》使前一线索可追，本纪与编年材料主要确证政权年序和精英行动；对基层执行、疆界和人口数量的沉默不得以叙事补足；《水经注》；墓志、城址与十六国史辑佚限制后一线索的说法，僧传、道教文本、造像记与遗址的成书或营造年代须分开，个案不能量化信仰或政策落实。

沿着公元363年（年序桥接），东晋、前燕与前秦的\eventanchor{JH-DENS-363-A}{前燕、前秦与东晋围绕河北、关中和淮汉走廊轮番调兵，扩张依赖州郡供给与地方合作}指向前燕、前秦与东晋围绕河北、关中和淮汉走廊轮番调兵，扩张依赖州郡供给与地方合作，\eventanchor{JH-DENS-363-B}{东晋、前燕与前秦的关隘、渡口、军镇和水陆道路的可通行性继续约束使节、军队和物资的行动}则把东晋、前燕与前秦的关隘、渡口、军镇和水陆道路的可通行性继续约束使节、军队和物资的行动接入同一条年序。 如此处理不是用空白填满史实，而是避免后来的转折抹去本年仍在运作的条件。 前者依据《晋书》哀海西简文诸帝纪；《晋书》载记；《资治通鉴·晋纪》；本纪与编年材料主要确证政权年序和精英行动；对基层执行、疆界和人口数量的沉默不得以叙事补足。后者参照《水经注》；墓志、城址与十六国史辑佚；纪传中的行军路线、兵数与控制范围常被简化或夸饰，地理材料只能提供有限交叉。

\eventanchor{JH-DENS-364-A}{前燕、前秦与东晋围绕河北、关中和淮汉走廊轮番调兵，扩张依赖州郡供给与地方合作}和\eventanchor{JH-DENS-364-B}{东晋、前燕与前秦的朝廷和州郡围绕户籍、田租、军粮与转运的安排进入材料}使公元364年（年序桥接）的东晋、前燕与前秦既保留前燕、前秦与东晋围绕河北、关中和淮汉走廊轮番调兵，扩张依赖州郡供给与地方合作，也不略去东晋、前燕与前秦的朝廷和州郡围绕户籍、田租、军粮与转运的安排进入材料。 两条线索共同说明，政权变化之外，地方社会和空间条件仍在塑造行动者可以选择的办法。 《晋书》哀海西简文诸帝纪；《晋书》载记；《资治通鉴·晋纪》使前一线索可追，本纪与编年材料主要确证政权年序和精英行动；对基层执行、疆界和人口数量的沉默不得以叙事补足；《水经注》；墓志、城址与十六国史辑佚限制后一线索的说法，志书、诏令与本纪可显示制度设计或战争需求，不能直接换算赋额、仓储或实际负担。

在公元365年（年序桥接），\eventanchor{JH-DENS-365-A}{前燕、前秦与东晋围绕河北、关中和淮汉走廊轮番调兵，扩张依赖州郡供给与地方合作}所见的前燕、前秦与东晋围绕河北、关中和淮汉走廊轮番调兵，扩张依赖州郡供给与地方合作与\eventanchor{JH-DENS-365-B}{东晋、前燕与前秦的流民、侨置户、军镇居民或地方家族的安置与编户问题构成社会背景}所见的东晋、前燕与前秦的流民、侨置户、军镇居民或地方家族的安置与编户问题构成社会背景并行发生在东晋、前燕与前秦的历史空间。 年序因此不把大事件当作一切的终点，而让制度、交通和社会关系继续承受压力。 《晋书》哀海西简文诸帝纪；《晋书》载记；《资治通鉴·晋纪》只能确认前一观察的基本年序，本纪与编年材料主要确证政权年序和精英行动；对基层执行、疆界和人口数量的沉默不得以叙事补足；《水经注》；墓志、城址与十六国史辑佚为后一观察提供交叉线索，墓志、家族文本和侨置名目各自有选择性，不能据此统计迁徙规模或概括整体社会态度。

公元366年（年序桥接），东晋、前燕与前秦的\eventanchor{JH-DENS-366-A}{前燕、前秦与东晋围绕河北、关中和淮汉走廊轮番调兵，扩张依赖州郡供给与地方合作}把前燕、前秦与东晋围绕河北、关中和淮汉走廊轮番调兵，扩张依赖州郡供给与地方合作留在年序中；\eventanchor{JH-DENS-366-B}{东晋、前燕与前秦的寺院、僧侣、道士和施主网络在材料中连接都城、州郡与边地}则从东晋、前燕与前秦的寺院、僧侣、道士和施主网络在材料中连接都城、州郡与边地观察同一年的约束。 如此处理不是用空白填满史实，而是避免后来的转折抹去本年仍在运作的条件。 《晋书》哀海西简文诸帝纪；《晋书》载记；《资治通鉴·晋纪》只能确认前一观察的基本年序，本纪与编年材料主要确证政权年序和精英行动；对基层执行、疆界和人口数量的沉默不得以叙事补足；《水经注》；墓志、城址与十六国史辑佚为后一观察提供交叉线索，僧传、道教文本、造像记与遗址的成书或营造年代须分开，个案不能量化信仰或政策落实。

\eventanchor{JH-DENS-367-A}{前燕、前秦与东晋围绕河北、关中和淮汉走廊轮番调兵，扩张依赖州郡供给与地方合作}在公元367年（年序桥接）提示前燕、前秦与东晋围绕河北、关中和淮汉走廊轮番调兵，扩张依赖州郡供给与地方合作仍未结案；\eventanchor{JH-DENS-367-B}{东晋、前燕与前秦的关隘、渡口、军镇和水陆道路的可通行性继续约束使节、军队和物资的行动}以东晋、前燕与前秦的关隘、渡口、军镇和水陆道路的可通行性继续约束使节、军队和物资的行动说明东晋、前燕与前秦的选择还受具体条件限制。 两条线索共同说明，政权变化之外，地方社会和空间条件仍在塑造行动者可以选择的办法。 前者依据《晋书》哀海西简文诸帝纪；《晋书》载记；《资治通鉴·晋纪》；本纪与编年材料主要确证政权年序和精英行动；对基层执行、疆界和人口数量的沉默不得以叙事补足。后者参照《水经注》；墓志、城址与十六国史辑佚；纪传中的行军路线、兵数与控制范围常被简化或夸饰，地理材料只能提供有限交叉。

把公元368年（年序桥接）放回连续叙事，\eventanchor{JH-DENS-368-A}{前燕、前秦与东晋围绕河北、关中和淮汉走廊轮番调兵，扩张依赖州郡供给与地方合作}保留前燕、前秦与东晋围绕河北、关中和淮汉走廊轮番调兵，扩张依赖州郡供给与地方合作，\eventanchor{JH-DENS-368-B}{东晋、前燕与前秦的朝廷和州郡围绕户籍、田租、军粮与转运的安排进入材料}让东晋、前燕与前秦的朝廷和州郡围绕户籍、田租、军粮与转运的安排进入材料继续约束东晋、前燕与前秦的行动。 如此处理不是用空白填满史实，而是避免后来的转折抹去本年仍在运作的条件。 证据不许把这两个观察写成全域事实：《晋书》哀海西简文诸帝纪；《晋书》载记；《资治通鉴·晋纪》受本纪与编年材料主要确证政权年序和精英行动；对基层执行、疆界和人口数量的沉默不得以叙事补足。限制，《水经注》；墓志、城址与十六国史辑佚亦受志书、诏令与本纪可显示制度设计或战争需求，不能直接换算赋额、仓储或实际负担限制。

公元369年（年序桥接），\eventanchor{JH-DENS-369-A}{前燕、前秦与东晋围绕河北、关中和淮汉走廊轮番调兵，扩张依赖州郡供给与地方合作}和\eventanchor{JH-DENS-369-B}{东晋、前燕与前秦的流民、侨置户、军镇居民或地方家族的安置与编户问题构成社会背景}共同避免把东晋、前燕与前秦压缩为单一都城事件：前者关注前燕、前秦与东晋围绕河北、关中和淮汉走廊轮番调兵，扩张依赖州郡供给与地方合作，后者关注东晋、前燕与前秦的流民、侨置户、军镇居民或地方家族的安置与编户问题构成社会背景。 这里保留的是可核的连续性，而非对基层反应、疆界或人数的无据补写。 证据不许把这两个观察写成全域事实：《晋书》哀海西简文诸帝纪；《晋书》载记；《资治通鉴·晋纪》受本纪与编年材料主要确证政权年序和精英行动；对基层执行、疆界和人口数量的沉默不得以叙事补足。限制，《水经注》；墓志、城址与十六国史辑佚亦受墓志、家族文本和侨置名目各自有选择性，不能据此统计迁徙规模或概括整体社会态度限制。


《晋书》载记、《资治通鉴》与地理材料可定主要转折；十六国疆域和路线多有异文。


\subsection{370—379年：材料所见的连续压力}
前秦扩张前的关中、河洛与江南没有一条可放心使用的直线；使节、军队和粮食都受关隘、渡口与地方选择约束。

沿着公元370年（年序桥接），东晋、前燕与前秦的\eventanchor{JH-DENS-370-A}{前燕、前秦与东晋围绕河北、关中和淮汉走廊轮番调兵，扩张依赖州郡供给与地方合作}指向前燕、前秦与东晋围绕河北、关中和淮汉走廊轮番调兵，扩张依赖州郡供给与地方合作，\eventanchor{JH-DENS-370-B}{东晋、前燕与前秦的寺院、僧侣、道士和施主网络在材料中连接都城、州郡与边地}则把东晋、前燕与前秦的寺院、僧侣、道士和施主网络在材料中连接都城、州郡与边地接入同一条年序。 年序因此不把大事件当作一切的终点，而让制度、交通和社会关系继续承受压力。 前者依据《晋书》哀海西简文诸帝纪；《晋书》载记；《资治通鉴·晋纪》；本纪与编年材料主要确证政权年序和精英行动；对基层执行、疆界和人口数量的沉默不得以叙事补足。后者参照《水经注》；墓志、城址与十六国史辑佚；僧传、道教文本、造像记与遗址的成书或营造年代须分开，个案不能量化信仰或政策落实。

在公元371年（年序桥接）的并行行动者之间，\eventanchor{JH-DENS-371-A}{前燕、前秦与东晋围绕河北、关中和淮汉走廊轮番调兵，扩张依赖州郡供给与地方合作}将前燕、前秦与东晋围绕河北、关中和淮汉走廊轮番调兵，扩张依赖州郡供给与地方合作落到时间位置，\eventanchor{JH-DENS-371-B}{东晋、前燕与前秦的关隘、渡口、军镇和水陆道路的可通行性继续约束使节、军队和物资的行动}从东晋、前燕与前秦的关隘、渡口、军镇和水陆道路的可通行性继续约束使节、军队和物资的行动补足资源与空间的视角。 它们将政治时间同资源和地方网络并置，却不把并置误作已被证实的直接因果。 材料边界须分别保留：《晋书》哀海西简文诸帝纪；《晋书》载记；《资治通鉴·晋纪》与《水经注》；墓志、城址与十六国史辑佚所见不同，且本纪与编年材料主要确证政权年序和精英行动；对基层执行、疆界和人口数量的沉默不得以叙事补足；纪传中的行军路线、兵数与控制范围常被简化或夸饰，地理材料只能提供有限交叉。

公元372年（年序桥接），\eventanchor{JH-DENS-372-A}{前秦向淮汉推进，东晋以荆扬军府和江淮水陆通道维持防务，战争准备成为年序主轴}和\eventanchor{JH-DENS-372-B}{东晋与前秦的朝廷和州郡围绕户籍、田租、军粮与转运的安排进入材料}共同避免把东晋与前秦压缩为单一都城事件：前者关注前秦向淮汉推进，东晋以荆扬军府和江淮水陆通道维持防务，战争准备成为年序主轴，后者关注东晋与前秦的朝廷和州郡围绕户籍、田租、军粮与转运的安排进入材料。 这里保留的是可核的连续性，而非对基层反应、疆界或人数的无据补写。 材料边界须分别保留：《晋书·孝武帝纪》；《晋书·苻坚载记》；《资治通鉴·晋纪》与《水经注》；城址、墓志与军事地理研究所见不同，且本纪与编年材料主要确证政权年序和精英行动；对基层执行、疆界和人口数量的沉默不得以叙事补足；志书、诏令与本纪可显示制度设计或战争需求，不能直接换算赋额、仓储或实际负担。

把公元373年（年序桥接）放回连续叙事，\eventanchor{JH-DENS-373-A}{前秦向淮汉推进，东晋以荆扬军府和江淮水陆通道维持防务，战争准备成为年序主轴}保留前秦向淮汉推进，东晋以荆扬军府和江淮水陆通道维持防务，战争准备成为年序主轴，\eventanchor{JH-DENS-373-B}{东晋与前秦的流民、侨置户、军镇居民或地方家族的安置与编户问题构成社会背景}让东晋与前秦的流民、侨置户、军镇居民或地方家族的安置与编户问题构成社会背景继续约束东晋与前秦的行动。 这使读者能把下一年的军事、行政或社会变化放回尚未消散的限制中。 就可说范围而言，前一判断止于《晋书·孝武帝纪》；《晋书·苻坚载记》；《资治通鉴·晋纪》，本纪与编年材料主要确证政权年序和精英行动；对基层执行、疆界和人口数量的沉默不得以叙事补足；后一判断止于《水经注》；城址、墓志与军事地理研究，墓志、家族文本和侨置名目各自有选择性，不能据此统计迁徙规模或概括整体社会态度。

\eventanchor{JH-DENS-374-A}{前秦向淮汉推进，东晋以荆扬军府和江淮水陆通道维持防务，战争准备成为年序主轴}在公元374年（年序桥接）提示前秦向淮汉推进，东晋以荆扬军府和江淮水陆通道维持防务，战争准备成为年序主轴仍未结案；\eventanchor{JH-DENS-374-B}{东晋与前秦的寺院、僧侣、道士和施主网络在材料中连接都城、州郡与边地}以东晋与前秦的寺院、僧侣、道士和施主网络在材料中连接都城、州郡与边地说明东晋与前秦的选择还受具体条件限制。 这两项观察不宣称另有一场未载战争或政策，只把已有压力的时间位置和可见面向分开。 此处的证据支点是《晋书·孝武帝纪》；《晋书·苻坚载记》；《资治通鉴·晋纪》和《水经注》；城址、墓志与军事地理研究；本纪与编年材料主要确证政权年序和精英行动；对基层执行、疆界和人口数量的沉默不得以叙事补足。，同时僧传、道教文本、造像记与遗址的成书或营造年代须分开，个案不能量化信仰或政策落实。

公元375年（年序桥接），东晋与前秦的\eventanchor{JH-DENS-375-A}{前秦向淮汉推进，东晋以荆扬军府和江淮水陆通道维持防务，战争准备成为年序主轴}把前秦向淮汉推进，东晋以荆扬军府和江淮水陆通道维持防务，战争准备成为年序主轴留在年序中；\eventanchor{JH-DENS-375-B}{东晋与前秦的关隘、渡口、军镇和水陆道路的可通行性继续约束使节、军队和物资的行动}则从东晋与前秦的关隘、渡口、军镇和水陆道路的可通行性继续约束使节、军队和物资的行动观察同一年的约束。 年序因此不把大事件当作一切的终点，而让制度、交通和社会关系继续承受压力。 材料边界须分别保留：《晋书·孝武帝纪》；《晋书·苻坚载记》；《资治通鉴·晋纪》与《水经注》；城址、墓志与军事地理研究所见不同，且本纪与编年材料主要确证政权年序和精英行动；对基层执行、疆界和人口数量的沉默不得以叙事补足；纪传中的行军路线、兵数与控制范围常被简化或夸饰，地理材料只能提供有限交叉。

在公元376年（年序桥接），\eventanchor{JH-DENS-376-A}{前秦向淮汉推进，东晋以荆扬军府和江淮水陆通道维持防务，战争准备成为年序主轴}所见的前秦向淮汉推进，东晋以荆扬军府和江淮水陆通道维持防务，战争准备成为年序主轴与\eventanchor{JH-DENS-376-B}{东晋与前秦的朝廷和州郡围绕户籍、田租、军粮与转运的安排进入材料}所见的东晋与前秦的朝廷和州郡围绕户籍、田租、军粮与转运的安排进入材料并行发生在东晋与前秦的历史空间。 它们不是由年份自动推出的因果，而是让前后事件之间的资源、道路、户籍或地方网络保持可见。 证据不许把这两个观察写成全域事实：《晋书·孝武帝纪》；《晋书·苻坚载记》；《资治通鉴·晋纪》受本纪与编年材料主要确证政权年序和精英行动；对基层执行、疆界和人口数量的沉默不得以叙事补足。限制，《水经注》；城址、墓志与军事地理研究亦受志书、诏令与本纪可显示制度设计或战争需求，不能直接换算赋额、仓储或实际负担限制。

\eventanchor{JH-DENS-377-A}{前秦向淮汉推进，东晋以荆扬军府和江淮水陆通道维持防务，战争准备成为年序主轴}和\eventanchor{JH-DENS-377-B}{东晋与前秦的流民、侨置户、军镇居民或地方家族的安置与编户问题构成社会背景}使公元377年（年序桥接）的东晋与前秦既保留前秦向淮汉推进，东晋以荆扬军府和江淮水陆通道维持防务，战争准备成为年序主轴，也不略去东晋与前秦的流民、侨置户、军镇居民或地方家族的安置与编户问题构成社会背景。 年序因此不把大事件当作一切的终点，而让制度、交通和社会关系继续承受压力。 就可说范围而言，前一判断止于《晋书·孝武帝纪》；《晋书·苻坚载记》；《资治通鉴·晋纪》，本纪与编年材料主要确证政权年序和精英行动；对基层执行、疆界和人口数量的沉默不得以叙事补足；后一判断止于《水经注》；城址、墓志与军事地理研究，墓志、家族文本和侨置名目各自有选择性，不能据此统计迁徙规模或概括整体社会态度。

沿着公元378年（年序桥接），东晋与前秦的\eventanchor{JH-DENS-378-A}{前秦向淮汉推进，东晋以荆扬军府和江淮水陆通道维持防务，战争准备成为年序主轴}指向前秦向淮汉推进，东晋以荆扬军府和江淮水陆通道维持防务，战争准备成为年序主轴，\eventanchor{JH-DENS-378-B}{东晋与前秦的寺院、僧侣、道士和施主网络在材料中连接都城、州郡与边地}则把东晋与前秦的寺院、僧侣、道士和施主网络在材料中连接都城、州郡与边地接入同一条年序。 这使读者能把下一年的军事、行政或社会变化放回尚未消散的限制中。 前者依据《晋书·孝武帝纪》；《晋书·苻坚载记》；《资治通鉴·晋纪》；本纪与编年材料主要确证政权年序和精英行动；对基层执行、疆界和人口数量的沉默不得以叙事补足。后者参照《水经注》；城址、墓志与军事地理研究；僧传、道教文本、造像记与遗址的成书或营造年代须分开，个案不能量化信仰或政策落实。

在公元379年（年序桥接）的并行行动者之间，\eventanchor{JH-DENS-379-A}{前秦向淮汉推进，东晋以荆扬军府和江淮水陆通道维持防务，战争准备成为年序主轴}将前秦向淮汉推进，东晋以荆扬军府和江淮水陆通道维持防务，战争准备成为年序主轴落到时间位置，\eventanchor{JH-DENS-379-B}{东晋与前秦的关隘、渡口、军镇和水陆道路的可通行性继续约束使节、军队和物资的行动}从东晋与前秦的关隘、渡口、军镇和水陆道路的可通行性继续约束使节、军队和物资的行动补足资源与空间的视角。 两条线索共同说明，政权变化之外，地方社会和空间条件仍在塑造行动者可以选择的办法。 证据不许把这两个观察写成全域事实：《晋书·孝武帝纪》；《晋书·苻坚载记》；《资治通鉴·晋纪》受本纪与编年材料主要确证政权年序和精英行动；对基层执行、疆界和人口数量的沉默不得以叙事补足。限制，《水经注》；城址、墓志与军事地理研究亦受纪传中的行军路线、兵数与控制范围常被简化或夸饰，地理材料只能提供有限交叉限制。


《晋书》、载记与《资治通鉴》是骨架；墓志和宗教材料只能补出局部。
